\documentclass[12pt]{article}

% ---- Formatting per assignment instructions ----
% Times New Roman, 12pt, 0.8" margins, 1.5 line spacing, 2-3 pages
\usepackage[margin=0.8in]{geometry}
\usepackage{mathptmx} % Times New Roman-compatible font
\usepackage{setspace}
\onehalfspacing

\usepackage{hyperref}
\usepackage{graphicx}
\usepackage{enumitem}

\title{\vspace{-2em}CS/SE 4AL3: Project Proposal}
\author{}
\date{}

\begin{document}

% ---- TA consulted (required at top of document) ----
\noindent \textbf{TA(s) Consulted:} [Name(s) of TA(s) you spoke to]

\vspace{0.5em}
\begin{center}
{\LARGE \textbf{Project Title Here}}
\end{center}

% ============================================================
\section{Team Members}
\begin{itemize}[noitemsep]
    \item Full Name 1 -- Student ID -- email@mcmaster.ca
    \item Full Name 2 -- Student ID -- email@mcmaster.ca
    \item Full Name 3 -- Student ID -- email@mcmaster.ca
\end{itemize}

% ============================================================
\section{Task Overview}
% Item 2: Task title and overview, significance, what makes it challenging
[Describe the project at a high level. State the task title, what problem you are solving,
why it matters, and what makes it non-trivial/challenging.]

% ============================================================
\section{Task Definition}
% Item 3: type of data, classification/regression/generation, number of classes, single/multi-label
[State clearly: What type of data are you using (text, image, tabular, audio, etc.)?
Is this a classification, regression, or generation task? If classification, how many classes,
and is it single-label or multi-label?]

% ============================================================
\section{Problem Description and Impact}
% Item 4: Describe the problem, its impact on the real world, and why it is challenging
[Explain the problem in more depth, its real-world significance/impact, and elaborate on
the technical or practical challenges involved.]

% ============================================================
\section{Data Source(s) and Collection Plan}
% Item 5
[Describe where your data comes from. Include:
\begin{itemize}[noitemsep]
    \item Links to the data source(s)
    \item If scraped: method and how you will respect terms-of-service / robots.txt
    \item If via API: how you will handle authentication and rate limiting
    \item If from Kaggle: BOTH the Kaggle link AND the original source link
    \item Available metadata for the corpus
    \item If data is unlabeled: your plan for labeling it, and estimated time per instance
    \item If not collected by you: details on how it was originally collected/annotated/preprocessed
    \item Whether you will use the full dataset or a subset, and your justification
\end{itemize}

% ============================================================
\section{Dataset Size and Example Data Points}
% Item 6: expected size (>= 1k data points) and 3 example data points with labels
[State the expected number of data points in your dataset (minimum 1,000).]

\subsection*{Example Data Points}
\begin{enumerate}[noitemsep]
    \item \textbf{Input:} [example 1] \quad \textbf{Label:} [label 1]
    \item \textbf{Input:} [example 2] \quad \textbf{Label:} [label 2]
    \item \textbf{Input:} [example 3] \quad \textbf{Label:} [label 3]
\end{enumerate}

% ============================================================
\section{Proposed Solution}
% Item 7: modeling approach, features/labels, candidate models, existing solutions,
% 2-5 cited sources, evaluation plan, libraries
[Describe your planned approach:
\begin{itemize}[noitemsep]
    \item What features and target labels will you use?
    \item What model(s) will you try (no simple linear regression)?
    \item What existing solutions/approaches exist for this problem?
    \item How will you evaluate whether your model is good (metrics, validation strategy)?
    \item What libraries/frameworks do you intend to use?
\end{itemize}

\subsection*{Sources of Inspiration}
% Must cite 2-5 sources (papers, books, ML challenges)
\begin{enumerate}[noitemsep]
    \item Author(s), ``Title,'' Venue/Year.
    \item Author(s), ``Title,'' Venue/Year.
\end{enumerate}

% ============================================================
\section{TA Consultation Summary}
% Item 8: paragraph on what you discussed with your TA, their advice, and how it shaped your plans
[Summarize what you discussed with your TA, the advice/feedback they provided, and how it
influenced or changed your project plans.]

% ============================================================
% Note: The Team Contract (Item 9) must be submitted as a SEPARATE PDF, not part of this document.
% See the assignment instructions, Section 4, for required contents and signatures.

\end{document}
